\chapter*{Glossary}
\addcontentsline{toc}{chapter}{GLOSSARY} 

\textbf{Microcontroller:}  
A compact integrated circuit designed to govern a specific operation in an embedded system.

\textbf{Embedded System:}  
A combination of hardware and software designed to perform a dedicated function within a larger system.

\textbf{Seven-Segment Display:}  
An electronic display device used to display decimal numerals using seven individual LED segments.

\textbf{Multiplexing:}  
A technique that allows multiple display units to share common data lines by activating them sequentially.

\textbf{Hardware Timer:}  
A built-in peripheral in a microcontroller used to measure time intervals and generate delays.

\textbf{Interrupt:}  
A signal that temporarily halts the main program execution to execute a specific service routine.

\textbf{CTC Mode:}  
Clear Timer on Compare mode used to reset the timer counter when a predefined value is reached.

\textbf{Internal RC Oscillator:}  
An on-chip clock source that generates timing signals without the need for external components.

\textbf{Proteus:}  
A simulation software used for designing and testing electronic circuits and embedded systems.

\textbf{Microchip Studio:}  
An integrated development environment for developing and debugging AVR microcontroller applications.
